\section{Governance and Reproducibility}
\label{sec:governance}

For a small, single-maintainer benchmark the two credibility risks that matter most
are silent judge drift and run cherry-picking: re-running until the numbers look
good, or quietly swapping a judge model out from under a published ranking. The
governance policy exists to make both visible and to make corrections auditable.
Each policy below is enforced in code where stated.

\subsection{Pre-registration}

Before the first agent, simulator, or judge call, the run writes a
\texttt{pre\_registration.json} fixing the run's definition: the run identity and a
UTC timestamp, the requested model roster (name, model id, provider), the scenario
set (domains, per-domain scenario IDs, and a corpus-level SHA-256 over the
canonical serialized scenario set, with a recomputable per-scenario index), the
configured judge panel, the reliability-run count, the seeds and temperatures, and
the judge-prompt mode. A test asserts the file exists on disk at the moment the
first model is dispatched, so the ordering is enforced, not merely intended. Because
the definition is fixed before any number is known, the maintainer cannot
retroactively choose which run counts. The record states explicitly that the user
and tool simulators run unseeded and the user simulator at temperature above zero,
so runs are not bit-for-bit reproducible; the honest caveat is part of the artifact.

\subsection{Judge pinning}

Judge versions are pinned per run, and the model a provider \emph{actually served}
is recorded, not just the model requested. Hosted model IDs (OpenRouter slugs in
particular) can be silently re-pointed to a different snapshot or quantization
without the calling code changing, so a pinned request ID alone would not catch the
drift. Every judge call serializes a \texttt{resolved\_model} into the artifact,
and the agent under test records its resolved model too. A judge version change is
treated as a methodology change (Section~\ref{sec:versioning}).

\subsection{No silent retraction or rerun}

A published run is never deleted and never quietly re-run. Corrections are
append-only: a problem with a published run is fixed by computing a new, dated run
and recording the reason in a changelog, not by overwriting the original. The old
numbers stay visible alongside the correction. This is the defense against the
re-run-until-it-looks-right failure mode, because there is no clean slate to re-run
toward.

\subsection{Contamination and the private holdout}

Scenario authors are never contestants, enforced by the family-aware guard of
Section~\ref{sec:design}. Authorship records are never altered. Synthetic scenarios
are a strong native contamination defense, but the public corpus is on GitHub, so a
future model could in principle train on it. The committed defense is a private
holdout: a fresh scenario subset, authored to the same v0.2 schema and quality bar,
stored outside the repository and never published. It is pulled in at evaluation
time via \texttt{--holdout-dir} and graded by the same simulator, judges, and state
checks as the public corpus. When a holdout is run it gets its own
pre-registration entry carrying the corpus SHA-256 \emph{and count only}, no
scenario IDs and no per-scenario index, so the set is pinned and tamper-evident
without revealing its content. Holdout rows are split out before aggregation, so the
public efficacy and CLEAR rankings are computed over the public corpus only; the
board publishes per model the public score, the holdout score, and the gap
($\text{gap} = \text{public} - \text{holdout}$). A materially positive gap is the
overfitting tripwire. No holdout scenario ID, text, ground truth, or per-scenario
score reaches any published surface; the public CI refuses to start if a holdout
directory is configured, because its artifact upload would make transcripts
downloadable, so holdout runs are local-only by construction. This is the
SWE-bench-Pro held-out-split blueprint~\cite{swebenchpro2025} in miniature.

\subsection{Versioning}

\label{sec:versioning}
A benchmark version is the unit of comparability; scores are comparable only within
a version. A version bump is required for any corpus change (adding, removing, or
modifying scenarios; a published scenario is immutable, the MTEB
lesson~\cite{mteb2025}), any judge change (a different judge is a different measuring
instrument), or any rubric change (rubrics, dimension weights, or efficacy
composition). Cross-version scores are not comparable, which is already structural
in the field-relative CLEAR normalization.

\subsection{What every run publishes}

A published score is only trustworthy if the evidence behind it is inspectable. Each
run persists: the pre-registration (above); per-evaluation artifacts for every
(scenario, model, run index) containing the full transcript in true conversational
order, every judge's output (name, rubric, score, reasoning, parse-failed flag, raw
parsed response, resolved model), simulation metadata including the
completion-decoupling fields, the deterministic state-check result, and the
scenario's domain, category, and holdout flag; an OpenInference-attributed
OpenTelemetry trace export written as JSONL, loadable into Arize Phoenix or any OTel
reader; and a run manifest linking back to the pre-registration by path and hash, so
the two are a verifiable pair. All four are uploaded as workflow artifacts on every
weekly run.

\subsection{Maintenance posture}

The benchmark runs on automated weekly GitHub Actions, so the commit timestamps are
the liveness signal rather than a manual quarterly re-run that depends on free time.
Scope is deliberately narrow (two domains for a one-person-plus-community operation,
not dozens of sub-industries), since scope creep is the primary cause of
solo-project abandonment. The intended long-term model is publish-the-harness,
others-run-it, submit-via-PR, with versioned releases as content events.
